%% kapa: exam=1 round=26 year=2015 period=1 form=A subject=civil_law subject_ko=민법 seq=23 number=23
%% source: Q-Net 감정평가사 기출문제 (cst00309) artlSeq=5138001 file=extracted/1008/5138001/제26회_감정평가사 1차_기출/제26회_감정평가사 1차_기출/1교시_감평_A형.pdf
%% conversion: kordoc 4.24.51 -> tools/gen1.py (auto); review=proofread-sonnet; flags=none
\begin{question}{23}
부동산 물권변동에 관한 설명으로 옳지 않은 것은? (다툼이 있으면 판례에 따름)
\begin{choices}
  \item 甲이 매매를 원인으로 하는 소유권이전등기소송에서 승소의 확정판결을 얻었더라도 이전등기 전에는 소유권을 취득하지 못한다.
  \item 재단법인의 설립을 위해 부동산을 출연한 경우, 출연자와 재단법인 사이에서도 그 부동산은 소유권 이전등기 없이는 재단법인의 소유가 되지 않는다.
  \item 甲이 건물을 신축한 후 乙에게 양도하고 乙 명의로 보존등기를 한 경우 乙은 건물의 소유권을 취득한다.
  \item 공용징수에 의한 부동산 소유권의 취득에는 등기를 요하지 않는다.
  \item 甲이 乙에게 부동산을 매도하고 소유권이전등기를 한 후 계약을 해제하였으나 그 말소등기 전에 乙이 선의의 丙에게 매도하고 이전등기한 경우, 甲은 丙에게 등기의 말소를 청구할 수 없다.
\end{choices}
\end{question}
\answer{2}
